% ======================================================================
% 12. přednáška – 19. 12. 2022
% Kruhové urychlovače (pokračování): synchrotron, synchrotronové záření
% ======================================================================
\section{Kruhové urychlovače: synchrotron a synchrotronové záření}

\strana{P12-1}
\emph{12. přednáška, 19. prosince 2022.}

\subsection{Kruhové urychlovače -- pokračování}

Minule jsme probrali urychlovače, v nichž částice obíhají v konstantním
magnetickém poli (cyklotron) -- jejich dráha se s rostoucí energií
rozšiřuje a při relativistických energiích přestane souhlasit frekvence
oběhu s frekvencí urychlujícího pole. Pro nejvyšší energie potřebujeme jiný
princip: udržet částici na \emph{stálé} kruhové dráze a magnetické pole
zvyšovat spolu s její energií. To je \kw{synchrotron}.

\subsubsection*{Synchrotron}

Synchrotron je prstenec (vakuová trubice) s dipólovými magnety
(\(M\)), které zakřivují dráhu, a s urychlovacími dutinami
(\(RF\), radiofrekvenční rezonátory), v nichž částice při každém oběhu
získá malou porci energie:
\begin{center}
\begin{tikzpicture}[scale=1.0]
  \draw[thick] (0,0) ellipse (2.4 and 1.75);
  % RF dutina
  \begin{scope}[rotate around={0:(0,0)}]
    \draw[thick,fill=gray!15] (-0.35,1.6) rectangle (0.35,1.95);
    \node[above] at (0,1.95) {\(RF\)};
  \end{scope}
  % magnety
  \draw[thick,fill=Blue!15,rotate around={-28:(1.2,1.52)}]
     (0.95,1.35) rectangle (1.45,1.7);
  \node at (1.75,2.0) {\(M\)};
  \draw[thick,fill=Blue!15,rotate around={-60:(2.0,0.97)}]
     (1.75,0.8) rectangle (2.25,1.15);
  \node at (2.6,1.35) {\(M\)};
  \draw[-{Stealth[length=5pt]},thick,red!70!black] (-1.4,-1.42) arc[start angle=234,end angle=270,x radius=2.4,y radius=1.75];
\end{tikzpicture}
\end{center}
Pro částici s jednotkovým nábojem na kružnici o poloměru \(R\) v poli
\(B\) platí (rovnováha Lorentzovy a dostředivé síly, \(p=eBR\)) praktický
vztah
\[
  \vzorec{p=0{,}3\,B\cdot R}\qquad\Longrightarrow\qquad
  \frac{E}{B}=\text{konst},
\]
kde \(p\) je v \(\mathrm{GeV}/c\), \(B\) v teslech a \(R\) v metrech
(číslo \(0{,}3\) je \(c\) v jednotkách \(10^9\,\mathrm{m/s}\), tj.\
\(e\cdot c/(1\,\mathrm{GeV})\approx0{,}3\)). Protože \(R\) je pevné a pro
ultrarelativistické částice \(E\approx pc\), musí magnetické pole růst
\emph{úměrně energii} -- odtud podmínka \(E/B=\text{konst}\) (viz
Davídkovy slajdy „Urychlovače“). Současně se musí přizpůsobovat i frekvence
RF pole (dokud částice není ultrarelativistická a rychlost se ještě mění).

\begin{priklad}[Odhad pro LHC]
\(B\approx8{,}3\,\mathrm{T}\), poloměr zakřivení v dipólech
\(R\approx2{,}8\,\mathrm{km}\): \(p\approx0{,}3\cdot8{,}3\cdot2800\approx
7000\,\mathrm{GeV}/c\), tedy právě \(7\,\mathrm{TeV}\) na svazek.
\end{priklad}

\subsection{Synchrotronové záření}

Každý nabitý objekt, který se pohybuje se zrychlením, vyzařuje. Pro
kruhové urychlovače elektronů je to hlavní omezení dosažitelné energie --
energii vyzářenou při každém oběhu musí RF dutiny znovu dodat. Výchozím
bodem je relativistické zobecnění \kw{Larmorova vzorce}: vyzařovaný výkon
\[
  P_s=\frac{e^2}{6\pi\varepsilon_0 m^2c^3}\left(\dv{\vec p}{t}\right)^2
\]
(nerelativistický tvar pro malé rychlosti, kde \(\dv*{\vec p}{t}=m\vec a\)).
Aby byl vzorec relativisticky správný, nahradíme derivaci podle času
derivací podle \kw{vlastního času} \(\tau\), \(\dd\tau=\dd t/\gamma\), a
čtverec trojvektoru čtvercem čtyřvektoru hybnosti \(p_\mu=(E/c,\vec p)\):
\[
  \left(\dv{p_\mu}{\tau}\right)^2\;\longrightarrow\;
  \left(\dv{\vec p}{\tau}\right)^2-\frac{1}{c^2}\left(\dv{E}{\tau}\right)^2 .
\]
(Znaménko je zvoleno tak, aby výraz byl kladný pro prostorupodobné
čtyřzrychlení; v klidové soustavě přechází na původní
\((\dv*{\vec p}{t})^2\).) Dva důležité speciální případy:

\subsubsection*{a) Kruhový urychlovač}
Při oběhu se energie téměř nemění, \(\dv*{E}{t}\approx0\); mění se jen
\emph{směr} hybnosti. Pootočí-li se hybnost o úhel \(\dd\alpha\), změní se
o \(\dd p=p\cdot\dd\alpha\), takže
\[
  \dv{p}{t}=p\cdot\dv{\alpha}{t}=p\cdot\omega=p\,\frac{v}{R} .
\]
\strana{P12-2}
Dosadíme (\(\dd\tau=\dd t/\gamma\) dává faktor \(\gamma^2\)):
\[
  P_s=\frac{e^2c\gamma^2}{6\pi\varepsilon_0 m^2c^4}\,\frac{p^2v^2}{R^2}
  \doteq\frac{e^2\gamma^4c}{6\pi\varepsilon_0R^2},
\]
kde jsme použili \(p=\gamma mv\) a \(v\approx c\) (přesně by zde bylo
\(\gamma^4\beta^4\)). Energie ztracená za jeden oběh je výkon krát doba
oběhu \(2\pi R/c\):
\[
  \Delta E\big|_{1\,\text{oběh}}=\oint P_s\,\dd t
  =\frac{e^2\gamma^4c}{6\pi\varepsilon_0R^2}\,\frac{2\pi R}{c}
  =\vzorec{\frac{e^2\gamma^4}{3\varepsilon_0R}} .
\]
Klíčová je čtvrtá mocnina \(\gamma=E/(mc^2)\): při stejné energii vyzařuje
elektron \((m_p/m_e)^4\approx10^{13}\)krát více než proton. Proto jsou
synchrotronovým zářením limitovány hlavně elektronové urychlovače, zatímco
protonové (LHC) jsou limitovány silou magnetů. Ztrátu lze zmenšit jen
zvětšením poloměru \(R\) -- ten se ale uplatní pouze lineárně.

\begin{priklad}[LEP]
Pro elektrony je praktický tvar
\(\Delta E[\mathrm{keV}]\approx88{,}5\,E^4[\mathrm{GeV}]/R[\mathrm{m}]\).
\begin{itemize}
  \item LEP-1 (\(45\,\mathrm{GeV}\), \(R=3000\,\mathrm{m}\)):
    \(\Delta E\doteq120\,\mathrm{MeV}\) za oběh,
  \item LEP-2 (\(100\,\mathrm{GeV}\), \(R=3000\,\mathrm{m}\)):
    \(\Delta E\doteq2{,}95\,\mathrm{GeV}\) za oběh.
\end{itemize}
Kontrola: \(88{,}5\cdot45^4/3000\approx1{,}2\cdot10^5\,\mathrm{keV}\) a
\(88{,}5\cdot10^8/3000\approx2{,}95\cdot10^6\,\mathrm{keV}\). Zvýšení
energie jen \(2{,}2\)krát zvýšilo ztráty \(\approx25\)krát; při LEP-2 tak
elektron za každý oběh ztratil asi \(3\,\%\) své energie, což bylo
prakticky na hranici možností RF systému. Proto se další \(e^+e^-\)
urychlovače plánují jako lineární (nebo s mnohem větším obvodem, FCC-ee).
\end{priklad}

\subsubsection*{b) Lineární urychlovač}
Zde je síla rovnoběžná s rychlostí, \(\vec\beta\parallel\vec p\parallel
\dv*{\vec p}{t}\), a mění se velikost hybnosti i energie. Protože
\(\dd E=v\,\dd p\), je \((\dv*{E}{t})^2=c^2\beta^2(\dv*{p}{t})^2\), a tedy
\[
  P_s=\frac{e^2\gamma^2c}{6\pi\varepsilon_0m^2c^4}\left(\dv{p}{t}\right)^2
  \left(1-\oprava{\beta^2}{\frac{1}{\beta^2}}\right)
  =\frac{e^2c}{6\pi\varepsilon_0m^2c^4}\left(\dv{E}{x}\right)^2 .
\]
Faktor \(\gamma^2(1-\beta^2)=1\) se vykrátí a použili jsme
\(\dv*{p}{t}=\dv*{E}{x}\) (obojí je síla, \(\dd E=F\,\dd x\)). Výkon
vyzařovaný v lineárním urychlovači tedy \emph{nezávisí na energii}, jen na
urychlovacím gradientu \(\dv*{E}{x}\).
\pozor{Pozor:} se členem \((1-1/\beta^2)\), jak je napsán v poznámkách, by
výraz vyšel záporný a výsledek by neplatil; správně je \((1-\beta^2)\).

\begin{priklad}[ILC]
Plánovaný lineární urychlovač ILC má gradient \(30\,\mathrm{MV/m}\),
tj.\ \(\dv*{E}{x}=30\,\mathrm{MeV/m}\approx4{,}8\cdot10^{-12}\,\mathrm{N}\).
Dosazením pro elektron
\[
  P_s=\frac{(1{,}6\cdot10^{-19})^2\,(4{,}8\cdot10^{-12})^2}
  {6\pi\cdot8{,}85\cdot10^{-12}\cdot(9{,}11\cdot10^{-31})^2\cdot(3\cdot10^{8})^3}
  =1{,}6\times10^{-16}\,\mathrm{W}.
\]
Za dobu průletu urychlovačem (\(\sim10\,\mathrm{km}/c\approx
3\cdot10^{-5}\,\mathrm{s}\)) to je zcela zanedbatelných
\(\sim10^{-20}\,\mathrm{J}\approx0{,}03\,\mathrm{eV}\). V lineárních
urychlovačích synchrotronové záření nehraje roli.
\end{priklad}

\begin{vysvetleni}[Využití synchrotronového záření]
Co je pro částicové fyziky ztrátou, je pro jiné obory cenný zdroj:
synchrotronové záření je intenzivní, kolimované (do kužele o úhlu
\(\sim1/\gamma\)) a laditelné od infračervené oblasti po tvrdé rentgenové
záření. Specializované zdroje (ondulátory, wigglery) se používají
v materiálovém výzkumu, chemii a biologii (viz Davídkovy slajdy
„Urychlovače (7)“).
\end{vysvetleni}
